% Part: computability
% Chapter: recursive-functions
% Section: pr-functions-computable

\documentclass[../../../include/open-logic-section]{subfiles}

\begin{document}

\olfileid{cmp}{rec}{cmp}
\olsection{原始递归函数是可计算的}

假设函数 $h$ 由原始递归定义
\begin{eqnarray*}
h(\vec x, 0)     & = & f(\vec x) \\
h(\vec x, y + 1) & = & g(\vec x, y, h(\vec x, y))
\end{eqnarray*}
，且函数 $f$ 和 $g$ 是可计算的。（我们用 $\vec x$ 简写 $x_0$, \dots, $x_{k-1}$。）那么 $h(\vec x, 0)$ 显然可计算，因为它即为 $f(\vec x)$，而由假设 $f$ 可计算。$h(\vec x, 1)$ 亦可计算，由于 $1 = 0 + 1$，故 $h(\vec x, 1)$ 即为
\begin{align*}
 h(\vec x, 1) & = g(\vec x, 0, h(\vec x, 0)) =  g(\vec x, 0, f(\vec x)).
\intertext{我们可以照此继续计算}
h(\vec x, 2) & = g(\vec x, 1, h(\vec x, 1)) = g(\vec x, 1, g(\vec x, 0, f(\vec x)))\\
h(\vec x, 3) & = g(\vec x, 2, h(\vec x, 2)) = g(\vec x, 2, g(\vec x, 1, g(\vec x, 0, f(\vec x))))\\
h(\vec x, 4) & = g(\vec x, 3, h(\vec x, 3)) = g(\vec x, 3, g(\vec x, 2, g(\vec x, 1, g(\vec x, 0, f(\vec x)))))\\
& \vdots
\end{align*}
因此，一般地，要计算 $h(\vec x, y)$，只需依次计算 $h(\vec x, 0)$, $h(\vec x, 1)$, \dots，直至得到 $h(\vec x, y)$。

由此可见，若函数 $f$ 和 $g$ 可计算，则原始递归定义产生的新函数亦可计算。若函数 $f$ 和 $g_i$ 可计算，则函数的复合亦生成可计算函数。

由于基本函数 $\Zero$、$\Succ$ 和 $\Proj{n}{i}$ 都是可计算的，且复合与原始递归均能由可计算函数生成可计算函数，故每个原始递归函数都是可计算的。

\end{document}